% BHU PYQ -- Manually transcribed & error-corrected
\documentclass[11pt,a4paper]{article}

\usepackage[a4paper,top=2.5cm,bottom=2.5cm,left=2.8cm,right=2.8cm,headheight=14pt]{geometry}
\usepackage{amsmath,amssymb}
\usepackage[T1]{fontenc}
\usepackage[utf8]{inputenc}
\usepackage{lmodern}
\usepackage{microtype}
\usepackage{fancyhdr}
\usepackage{enumitem}
\usepackage{hyperref}
\usepackage{lastpage}
\usepackage{parskip}

\hypersetup{colorlinks=true,urlcolor=black,linkcolor=black,
  pdftitle={CHB-401: Inorganic & Organic Chemistry-II -- BHU PYQ},pdfauthor={Banaras Hindu University}}

\pagestyle{fancy}\fancyhf{}
\renewcommand{\headrulewidth}{0.4pt}\renewcommand{\footrulewidth}{0.4pt}
\fancyhead[L]{\small   \textit{Banaras Hindu University}}
\fancyhead[R]{\small   \textit{CHB-401: Inorganic \& Organic Chemistry-II}}
\fancyfoot[C]{\small Page  \thepage\ of \pageref{LastPage}}


\newlist{parts}{enumerate}{2}
\setlist[parts,1]{label=(\alph*),leftmargin=*,itemsep=5pt,topsep=3pt}
\setlist[parts,2]{label=(\roman*),leftmargin=*,itemsep=3pt,topsep=2pt}

\newcommand{\pts}[1]{\hfill   \textit{[#1]}}


\begin{document}


\begin{center}
  {\large\bfseries BANARAS HINDU UNIVERSITY}\\[2pt]
  {
\normalsize B.Sc. (Hons.) Semester IV Examination, 2013-14}\\[6pt]
  \rule{0.8\linewidth}{0.5pt}\\[6pt]
  {\large\bfseries Paper No. CHB-401: Inorganic \& Organic Chemistry-II}\\[4pt]
  {
\normalsize Time: 3 Hours\hfill Maximum Marks: 70}
\end{center}
\rule{\linewidth}{0.4pt}
\small



\noindent   \textit{Instructions: Answer five questions in total, including Question~1 from each section which is compulsory. Use separate answer books for Section~A and Section~B.}

\normalsize
\bigskip


\begin{center}
  {\large\bfseries SECTION A}\\[2pt]
  {
\normalsize (Inorganic Chemistry-II -- Marks: 35)}
\end{center}
\rule{0.4\linewidth}{0.2pt}
\smallskip




\noindent   \textbf{Question 1.} Answer all parts of the following: \pts{5 $ \times$ 2 = 10}

\begin{parts}
  \item Why are alkali metals dissolved in liquid ammonia deep blue in color and good reducing agents?
  \item Why are $ \text{Ti}^{3+}$ complexes colored, whereas $ \text{Ti}^{4+}$ complexes are colorless?
  \item Why is the decrease in ionic size from $ \text{Ce}^{3+}$ to $ \text{Lu}^{3+}$ (lanthanide contraction) very small between adjacent elements?
  \item Which one is a stronger Lewis acid: $ \text{BCl}_3$ or $ \text{B(CH}_3 \text{)}_3$? Why?
  \item Draw the structures of linkage isomers of $[ \text{Co(NO}_2 \text{)(NH}_3 \text{)}_5]^{2+}$.
\end{parts}

\medskip\rule{\linewidth}{0.2pt}




\noindent   \textbf{Question 2.}

\begin{parts}
  \item State the Br\o nsted-Lowry and Lewis concepts of acids and bases. Identify the conjugate acid-base pairs in:
    \[  \text{NH}_3 +  \text{H}_2 \text{O} \rightleftharpoons  \text{NH}_4^+ +  \text{OH}^- \] \pts{6}
  \item In terms of the HSAB principle, which end of the $ \text{NCO}^-$ (cyanate) ion would you expect to coordinate to $ \text{Co}^{2+}$ (borderline acid) and $ \text{Ti}^{4+}$ (hard acid)? Explain. \pts{6}
\end{parts}

\medskip\rule{\linewidth}{0.2pt}




\noindent   \textbf{Question 3.}

\begin{parts}
  \item Why do transition metals exhibit variable oxidation states? Which oxidation states are exhibited by $ \text{Fe}$ and $ \text{Cr}$, and which ones are the most common? \pts{6}
  \item Write the IUPAC names of the following coordination compounds: \pts{6}
  
\begin{parts}
    \item $[ \text{Co(ONO)(NH}_3 \text{)}_5]^{2+}$
    \item $[( \text{NH}_3)_5 \text{Cr-OH-Cr(NH}_3)_5] \text{Cl}_5$
    \item $ \text{K}_2[ \text{Zn(NCS)}_4]$
  \end{parts}
\end{parts}

\medskip\rule{\linewidth}{0.2pt}




\noindent   \textbf{Question 4.}

\begin{parts}
  \item Compare and contrast the precipitation reactions in water and liquid ammonia solvents, giving two examples of each. \pts{6}
  \item Explain the important consequences of lanthanide contraction. Write the outer electronic configurations of $ \text{Ho}$ and $ \text{Lu}$ in the ground state. \pts{6}
\end{parts}


\newpage

\begin{center}
  {\large\bfseries SECTION B}\\[2pt]
  {
\normalsize (Organic Chemistry-III -- Marks: 35)}
\end{center}
\rule{0.4\linewidth}{0.2pt}
\smallskip




\noindent   \textbf{Question 5.} Answer all parts of the following: \pts{5 $ \times$ 2 = 10}

\begin{parts}
  \item Why is pyrrole more reactive than benzene towards electrophilic substitution reactions?
  \item Complete the following reaction:
    \[  \text{Pyridine} \xrightarrow{ \text{NaNH}_2} ? \]
  \item Complete the reaction and write down the product:
    \[  \text{Pyrrole} \xrightarrow{ \text{HCN, HCl}} ? \]
  \item How will you convert D-glucose into D-fructose?
  \item Define epimers and anomers of carbohydrates with examples.
\end{parts}

\medskip\rule{\linewidth}{0.2pt}




\noindent   \textbf{Question 6.}

\begin{parts}
  \item Write the name of the following reaction and explain the formation of the products with a suitable mechanism: \pts{6}
    \[  \text{Pyrrole} \xrightarrow{ \text{CHCl}_3,  \text{KOH}} ? \]
  \item Complete the following reaction, showing the stepwise mechanism: \pts{6}
    \[  \text{Naphthalene} \xrightarrow{ \text{Ac}_2 \text{O, AlCl}_3} ? \]
\end{parts}

\medskip\rule{\linewidth}{0.2pt}




\noindent   \textbf{Question 7.}

\begin{parts}
  \item Describe the Hantzsch pyridine synthesis with a detailed mechanism. \pts{6}
  \item Elucidate the structure of indigotin. How is it synthesized from anthranilic acid? \pts{6}
\end{parts}

\medskip\rule{\linewidth}{0.2pt}




\noindent   \textbf{Question 8.}

\begin{parts}
  \item How will you carry out the following conversions? \pts{6}
  
\begin{parts}
    \item Aldohexose to next higher aldohexose (Kiliani-Fischer synthesis)
    \item Aldohexose to next lower aldopentose (Wohl degradation)
  \end{parts}
  \item Compare the electrophilic aromatic substitution reactions of furan, pyrrole, and thiophene. \pts{6}
\end{parts}

\vfill
\rule{\linewidth}{0.4pt}

\begin{center}\small
  \href{https://bhu-exam.vercel.app/}{Click here to visit our website}\\[4pt]
  \href{https://me-aryan.vercel.app/}{Click here to visit our contributor}\\[4pt]
  \href{https://quashan-me.vercel.app/}{Click here to visit our contributor}
\end{center}

\end{document}
